\section*{Chapitre \ref{ch:g7:heart-and-circulation} --- Le cœur et la circulation}

\begin{solution}{exo:g7:heart-and-circulation:1}
\omterm{def:g7:heart-and-circulation:rooms}{Oreillette} droite~: elle reçoit le \omterm{prop:g4:journey-of-food:absorb}{sang} qui revient du corps.
\omterm{def:g7:heart-and-circulation:rooms}{Ventricule} droit~: il le pompe vers les \omterm{def:g7:how-animals-breathe:four}{poumons}. \omterm{def:g7:heart-and-circulation:rooms}{Oreillette} gauche~:
elle reçoit le \omterm{prop:g4:journey-of-food:absorb}{sang} riche en \omterm{def:g7:what-is-respiration:respiration}{dioxygène} des \omterm{def:g7:how-animals-breathe:four}{poumons}. \omterm{def:g7:heart-and-circulation:rooms}{Ventricule}
gauche~: il le pompe dans tout le corps --- c'est la cavité à la
paroi la plus épaisse des quatre.
\end{solution}

\begin{solution}{exo:g7:heart-and-circulation:2}
Ce sont des portes à sens unique qui empêchent le \omterm{prop:g4:journey-of-food:absorb}{sang} de refluer.
Le poum-tac, ce sont leurs fermetures par paires~: les \omterm{def:g7:heart-and-circulation:rooms}{valves}
d'entrée qui claquent quand les \omterm{def:g7:heart-and-circulation:rooms}{ventricules} chassent (poum), les
\omterm{def:g7:heart-and-circulation:rooms}{valves} de sortie qui claquent quand ils se relâchent (tac).
\end{solution}

\begin{solution}{exo:g7:heart-and-circulation:3}
Circuit des \omterm{def:g7:how-animals-breathe:four}{poumons}~: du \omterm{def:g7:heart-and-circulation:rooms}{ventricule} droit aux \omterm{def:g7:heart-and-circulation:vessels}{capillaires} des
\omterm{def:g7:how-animals-breathe:four}{poumons} (charger du \omterm{def:g7:what-is-respiration:respiration}{dioxygène}, décharger du \omterm{def:g7:what-is-respiration:respiration}{dioxyde de carbone}),
avec retour à l'\omterm{def:g7:heart-and-circulation:rooms}{oreillette} gauche. Circuit du corps~: du \omterm{def:g7:heart-and-circulation:rooms}{ventricule}
gauche aux \omterm{def:g7:heart-and-circulation:vessels}{capillaires} de tout le corps (livrer, collecter), avec
retour à l'\omterm{def:g7:heart-and-circulation:rooms}{oreillette} droite.
\end{solution}

\begin{solution}{exo:g7:heart-and-circulation:4}
\omterm{def:g7:heart-and-circulation:vessels}{Artères}~: des parois épaisses et élastiques --- elles emportent le
\omterm{prop:g4:journey-of-food:absorb}{sang} loin du \omterm{def:g4:heart-and-blood:heart}{cœur} sous pression. \omterm{def:g7:heart-and-circulation:vessels}{Capillaires}~: des parois d'une
seule \omterm{def:g5:first-look-at-cells:cell}{cellule} d'épaisseur --- le lieu de tous les échanges.
\omterm{def:g7:heart-and-circulation:vessels}{Veines}~: larges, à parois plus minces, avec des \omterm{def:g7:heart-and-circulation:rooms}{valves} à sens unique
--- elles ramènent le \omterm{prop:g4:journey-of-food:absorb}{sang} au \omterm{def:g4:heart-and-blood:heart}{cœur} à basse pression.
\end{solution}

\begin{solution}{exo:g7:heart-and-circulation:5}
Il anime le circuit du corps --- la boucle bien plus grande, jusqu'aux
orteils et au \omterm{def:g5:brain-in-command:system}{cerveau} et retour --- si bien qu'il doit pousser le
plus fort, et sa paroi est bâtie en conséquence.
\end{solution}

\begin{solution}{exo:g7:heart-and-circulation:6}
Dans les \omterm{def:g7:heart-and-circulation:vessels}{capillaires}. \omterm{def:g7:heart-and-circulation:vessels}{Capillaires} des \omterm{def:g7:how-animals-breathe:four}{poumons}~: \omterm{def:g7:what-is-respiration:respiration}{dioxygène} qui entre,
\omterm{def:g7:what-is-respiration:respiration}{dioxyde de carbone} qui sort~; ceux des \omterm{prop:g7:digestion-and-nutrients:absorption}{villosités}~: des \omterm{def:g7:digestion-and-nutrients:families}{nutriments}
qui entrent~; ceux des \omterm{def:g3:bones-and-muscles:muscle}{muscles}~: carburant et \omterm{def:g7:what-is-respiration:respiration}{dioxygène} débarqués,
\omterm{def:g7:what-is-respiration:respiration}{dioxyde de carbone} embarqué (ceux de la \omterm{def:g1:the-five-senses:sense}{peau}~: de la chaleur qui
sort).
\end{solution}

\begin{solution}{exo:g7:heart-and-circulation:7}
\omterm{def:g7:heart-and-circulation:rooms}{Ventricule} droit --- \omterm{def:g7:heart-and-circulation:vessels}{capillaires} des \omterm{def:g7:how-animals-breathe:four}{poumons} (\omterm{def:g7:what-is-respiration:respiration}{dioxygène} à bord) ---
\omterm{def:g7:heart-and-circulation:rooms}{oreillette} gauche --- \omterm{def:g7:heart-and-circulation:rooms}{ventricule} gauche --- \omterm{def:g7:heart-and-circulation:vessels}{artères} du cou --- les
\omterm{def:g7:heart-and-circulation:vessels}{capillaires} du \omterm{def:g5:brain-in-command:system}{cerveau} --- \omterm{def:g7:heart-and-circulation:vessels}{veines} du cou --- \omterm{def:g7:heart-and-circulation:rooms}{oreillette} droite.
L'échange du \omterm{def:g5:brain-in-command:system}{cerveau} prélève une part importante et constante de
\omterm{def:g7:what-is-respiration:respiration}{dioxygène} et de \omterm{prop:g7:organs-at-work:bill}{glucose} et laisse du \omterm{def:g7:what-is-respiration:respiration}{dioxyde de carbone} --- la
commande permanente d'un chef exigeant.
\end{solution}

\begin{solution}{exo:g7:heart-and-circulation:8}
Le \omterm{prop:g4:heart-and-blood:pulse}{pouls} est l'étirement élastique de la paroi de l'\omterm{def:g7:heart-and-circulation:vessels}{artère} à chaque
poussée du \omterm{def:g7:heart-and-circulation:rooms}{ventricule} --- une haute pression rencontrant une paroi
épaisse et souple. Arrivée aux \omterm{def:g7:heart-and-circulation:vessels}{veines}, la poussée s'est dépensée en
traversant les \omterm{def:g7:heart-and-circulation:vessels}{capillaires}~: une pression basse et régulière dans un
tube large et souple --- il ne reste rien pour battre.
\end{solution}

\begin{solution}{exo:g7:heart-and-circulation:9}
Le retour en montée~: depuis les \omterm{def:g1:my-body:parts}{jambes}, le \omterm{prop:g4:journey-of-food:absorb}{sang} veineux doit
grimper contre la pesanteur à basse pression, et les \omterm{def:g7:heart-and-circulation:rooms}{valves}
retiennent chaque gain. Chez le soldat parfaitement immobile, les
\omterm{def:g3:bones-and-muscles:muscle}{muscles} du mollet ne pressent jamais les \omterm{def:g7:heart-and-circulation:vessels}{veines}, la montée cale, le
\omterm{prop:g4:journey-of-food:absorb}{sang} stagne dans les \omterm{def:g1:my-body:parts}{jambes} et l'approvisionnement du \omterm{def:g5:brain-in-command:system}{cerveau} baisse
--- d'où l'évanouissement. Des pieds qui marchent sont des pompes à
\omterm{def:g7:heart-and-circulation:vessels}{veines}.
\end{solution}

\begin{solution}{exo:g7:heart-and-circulation:10}
Une \omterm{def:g7:heart-and-circulation:vessels}{artère} entrante venue du circuit du corps~; à l'intérieur, les
\omterm{def:g7:heart-and-circulation:vessels}{capillaires} des \omterm{prop:g7:digestion-and-nutrients:absorption}{villosités} absorbent les \omterm{def:g7:digestion-and-nutrients:families}{nutriments} du repas
(\cref{prop:g7:digestion-and-nutrients:absorption})~; la liaison
particulière~: sa \omterm{def:g7:heart-and-circulation:vessels}{veine} passe d'abord par le \omterm{prop:g7:digestion-and-nutrients:absorption}{foie} pour le tri et la
mise en banque, et c'est seulement ensuite que le \omterm{prop:g4:journey-of-food:absorb}{sang} rejoint le
retour général vers l'\omterm{def:g7:heart-and-circulation:rooms}{oreillette} droite. Le schéma raconte de
nouveau le chapitre de la \omterm{def:g4:journey-of-food:digestion}{digestion}.
\end{solution}

\begin{solution}{exo:g7:heart-and-circulation:11}
Le \omterm{prop:g4:journey-of-food:absorb}{sang} des cavités ne fait que traverser~; les parois épaisses du
\omterm{def:g4:heart-and-blood:heart}{cœur} lui-même, qui travaillent sans repos, ont besoin de leur propre
livraison \omterm{def:g7:heart-and-circulation:vessels}{capillaire} comme n'importe quel \omterm{def:g3:bones-and-muscles:muscle}{muscle} --- d'où de petites
\omterm{def:g7:heart-and-circulation:vessels}{artères} nourricières ramifiées sur le \omterm{def:g4:heart-and-blood:heart}{cœur}. Encrassées, elles
étranglent la conduite de carburant de la pompe~: le seul organe qui
ne peut jamais s'arrêter est celui que leur obstruction affame.
\end{solution}

\begin{solution}{exo:g7:heart-and-circulation:12}
Chacun des échanges dont vit le corps --- le \omterm{def:g7:what-is-respiration:respiration}{dioxygène} aux \omterm{def:g7:lungs-and-gas-exchange:alveoli}{alvéoles},
les \omterm{def:g7:digestion-and-nutrients:families}{nutriments} aux \omterm{prop:g7:digestion-and-nutrients:absorption}{villosités}, le carburant et les déchets à chaque
\omterm{def:g3:bones-and-muscles:muscle}{muscle} --- ne se fait qu'à travers les parois des \omterm{def:g7:heart-and-circulation:vessels}{capillaires}~; les
\omterm{def:g7:heart-and-circulation:vessels}{artères} et les \omterm{def:g7:heart-and-circulation:vessels}{veines} ne font qu'amener le \omterm{prop:g4:journey-of-food:absorb}{sang} à ces parois et
l'en éloigner, et le \omterm{def:g4:heart-and-blood:heart}{cœur} ne fait que le garder en mouvement. La
pompe et les autoroutes existent pour les voies fines comme des
cheveux~: arrête quelque part la circulation \omterm{def:g7:heart-and-circulation:vessels}{capillaire} et cet
organe meurt, si saine que soit la tuyauterie.
\end{solution}

\begin{solution}{pb:g7:heart-and-circulation:1}
\textbf{1.} Le \omterm{prop:g4:heart-and-blood:pulse}{pouls}~: le compte des poussées dans les \omterm{def:g7:heart-and-circulation:vessels}{artères} --- le
rythme de la pompe lu à un \omterm{def:g1:my-body:joint}{poignet}. La pression artérielle~: la
pression dans les \omterm{def:g7:heart-and-circulation:vessels}{artères} à la poussée et entre deux poussées --- la
charge des autoroutes. Le poum-tac du stéthoscope~: les \omterm{def:g7:heart-and-circulation:rooms}{valves}
entendues en train de se fermer --- l'état des portes.

\textbf{2.} Un \omterm{prop:g4:heart-and-blood:pulse}{pouls} de repos bas chez une patiente en forme signifie
un \omterm{def:g7:heart-and-circulation:rooms}{ventricule} puissant qui déplace plus de \omterm{prop:g4:journey-of-food:absorb}{sang} à chaque battement~:
le \omterm{def:g4:heart-and-blood:heart}{cœur} entraîné tourne calmement, avec de la réserve.

\textbf{3.} Les portes des \omterm{def:g7:heart-and-circulation:rooms}{valves}~: les \omterm{def:g7:heart-and-circulation:rooms}{valves} d'entrée à la poussée
des \omterm{def:g7:heart-and-circulation:rooms}{ventricules} (poum), les \omterm{def:g7:heart-and-circulation:rooms}{valves} de sortie à leur relâchement
(tac) --- quatre portes qui se ferment en deux paires, proprement.

\textbf{4.} \omterm{def:g7:heart-and-circulation:rooms}{Oreillette} droite, \omterm{def:g7:heart-and-circulation:rooms}{ventricule} droit, \omterm{def:g7:heart-and-circulation:vessels}{capillaires} des
\omterm{def:g7:how-animals-breathe:four}{poumons}, \omterm{def:g7:heart-and-circulation:rooms}{oreillette} gauche, \omterm{def:g7:heart-and-circulation:rooms}{ventricule} gauche, \omterm{def:g7:heart-and-circulation:vessels}{artères} du corps,
\omterm{def:g7:heart-and-circulation:vessels}{capillaires} du corps (par exemple d'un \omterm{def:g3:bones-and-muscles:muscle}{muscle}), \omterm{def:g7:heart-and-circulation:vessels}{veines} --- et retour
à l'\omterm{def:g7:heart-and-circulation:rooms}{oreillette} droite.

\textbf{5.} Chaque battement perd une partie de sa poussée à
contre-courant par la fuite, si bien que les livraisons du circuit
du corps diminuent à chaque battement~; le \omterm{def:g4:heart-and-blood:heart}{cœur} compense par plus de
battements et plus de travail, et l'endurance chute --- le souffle
est de la livraison perdue en reflux.

\textbf{6.} La pression artérielle monte (des autoroutes raides et
rétrécies résistent à la poussée) et l'étirement des \omterm{def:g7:heart-and-circulation:vessels}{artères} au
\omterm{prop:g4:heart-and-blood:pulse}{pouls} change~; le \omterm{def:g7:heart-and-circulation:rooms}{ventricule} gauche, qui pousse à chaque tournée
contre un circuit raidi, fait des heures supplémentaires et
s'épaissit.

\textbf{7.} L'encrassement des \omterm{def:g7:heart-and-circulation:vessels}{artères} nourricières du \omterm{def:g4:heart-and-blood:heart}{cœur} étrangle
la conduite de carburant du \omterm{def:g3:bones-and-muscles:muscle}{muscle} qui ne peut jamais se reposer ---
le chemin classique vers l'infarctus.

\textbf{8.} Les \omterm{def:g7:heart-and-circulation:rooms}{valves} à sens unique des \omterm{def:g7:heart-and-circulation:vessels}{veines} et la pompe
musculaire du mouvement, absente~: debout et immobile, les mollets
ne pressent jamais, les gains des \omterm{def:g7:heart-and-circulation:rooms}{valves} calent, et les \omterm{def:g1:my-body:joint}{chevilles} se
gorgent. La marche presse les \omterm{def:g7:heart-and-circulation:vessels}{veines} en rythme --- le soulagement
est mécanique.

\textbf{9.} Parce que le tour du corps dépense le \omterm{def:g7:what-is-respiration:respiration}{dioxygène} du
\omterm{prop:g4:journey-of-food:absorb}{sang}~: seul le circuit des \omterm{def:g7:how-animals-breathe:four}{poumons} peut le recharger. Sauter les
\omterm{def:g7:how-animals-breathe:four}{poumons} enverrait du \omterm{prop:g4:journey-of-food:absorb}{sang} épuisé en livraison avec une camionnette
vide --- la double boucle, c'est «~faire le plein puis livrer~»,
imposé par les deux côtés du \omterm{def:g4:heart-and-blood:heart}{cœur}.

\textbf{10.} \omterm{def:g7:heart-and-circulation:vessels}{Artère}~: épaisse, élastique, haute pression --- les
autoroutes vers l'extérieur. \omterm{def:g7:heart-and-circulation:vessels}{Capillaire}~: d'une \omterm{def:g5:first-look-at-cells:cell}{cellule} d'épaisseur,
se faufilant partout --- les voies de livraison où se font tous les
échanges. \omterm{def:g7:heart-and-circulation:vessels}{Veine}~: large, munie de \omterm{def:g7:heart-and-circulation:rooms}{valves}, basse pression --- le
retour.

\textbf{11.} L'ennui de chaque consultation comptait exactement là
où il touchait au service des \omterm{def:g7:heart-and-circulation:vessels}{capillaires}~: la fuite affamait les
livraisons, l'encrassement étranglait les conduites, la stagnation
bloquait les retours. Les lectures, les \omterm{def:g7:heart-and-circulation:rooms}{valves} et les autoroutes
sont des moyens~; la circulation \omterm{def:g7:heart-and-circulation:vessels}{capillaire} --- les échanges --- est
la fin.

\textbf{12.} Bouger chaque jour --- la pompe se renforce et les
autoroutes restent souples~; manger équilibré --- la paroi des
\omterm{def:g7:heart-and-circulation:vessels}{artères} reste nette de tout encrassement~; pas de fumée --- les
vaisseaux ne se rétrécissent ni ne se raidissent, et les places du
\omterm{prop:g4:journey-of-food:absorb}{sang} restent libres pour le \omterm{def:g7:what-is-respiration:respiration}{dioxygène}.
\end{solution}
